\documentclass[12pt, a4paper]{ctexart}
\usepackage{fontspec}
\usepackage{xcolor}
\usepackage{hyperref}
\usepackage{geometry}
\usepackage{enumitem}
\usepackage{parskip}
\usepackage{titlesec}
\usepackage{tcolorbox}
\tcbuselibrary{breakable}
\tcbset{autoparskip}

\usepackage{setspace}
\usepackage{listings}
\usepackage{amssymb}

\geometry{left=2.5cm, right=2.5cm, top=2.5cm, bottom=2.5cm}

\setmainfont{Times New Roman}
\setCJKmainfont{SimSun}

\hypersetup{
    colorlinks=true,
    linkcolor=blue!70!black,
    urlcolor=cyan,
    pdftitle={网络编程-逐题精讲解义},
    pdfauthor={专升本-fifine老师}
}

\definecolor{myblue}{RGB}{30,100,200}
\definecolor{lightblue}{RGB}{230,240,255}
\definecolor{darkblue}{RGB}{0,50,120}

\newtcolorbox{bluebox}[1][]{
    breakable,
    colback=lightblue,
    colframe=myblue,
    coltitle=darkblue,
    fonttitle=\bfseries,
    title={#1},
    boxrule=0.8pt,
    arc=4pt,
    left=6pt,
    right=6pt,
    top=6pt,
    bottom=6pt,
    before skip=8pt,
    after skip=8pt
}

\usepackage{etoolbox}
\BeforeBeginEnvironment{bluebox}{\par\vfill}%
\AfterEndEnvironment{bluebox}{\par\vfill}%

\titleformat{\section}
  {\normalfont\Large\bfseries\color{myblue}}{\thesection}{1em}{}
\titlespacing*{\section}{0pt}{3.5ex plus 1ex minus .2ex}{1.5ex plus .2ex}

\title{\textcolor{myblue}{\textbf{网络编程 - 逐题精讲解义}}}
\author{专升本-fifine老师 教学组}
\date{2026年03月05日}

\lstset{
    language=Java,
    basicstyle=\ttfamily\small,
    keywordstyle=\color{blue},
    commentstyle=\color{green!60!black},
    stringstyle=\color{orange},
    numbers=left,
    numberstyle=\tiny\color{gray},
    frame=single,
    breaklines=true,
    columns=fullflexible,
    showstringspaces=false
}

\newcommand{\code}[1]{\texttt{#1}}

\begin{document}

\maketitle
\thispagestyle{empty}
\newpage

\section{网络编程}

\begin{enumerate}[label=\textbf{\arabic*.}]

	\item \textbf{以下关于Java中Socket编程的说法错误的是}
	      \begin{enumerate}[label=\Alph*.]
		      \item 用于网络通信
		      \item ServerSocket用于服务器端
		      \item Socket用于客户端
		      \item 不需要处理异常
	      \end{enumerate}

	      \begin{bluebox}{答案与深度解析}
		      \textbf{答案：D. 不需要处理异常}

		      \textbf{【核心认知框架】}
		      \begin{itemize}[label={}, leftmargin=*, nosep]
			      \item \textbf{题目本质}：考查Socket编程的基本流程和异常处理
			      \item \textbf{关键区分点}：网络编程的不可靠性必须处理
			      \item \textbf{命题人意图}：测试对网络异常处理的重视程度
		      \end{itemize}

		      \textbf{【选项原理深度拆解】}

		      \begin{itemize}[leftmargin=*, nosep, label={}]
			      \item[A] \textbf{用于网络通信 — 正确}
			            \begin{itemize}[label={}, leftmargin=*, nosep]
				            \item \textbf{Socket概念}：网络通信端点的抽象（IP + Port）
				            \item \textbf{功能说明}：提供TCP/UDP网络编程能力
				            \item \textbf{使用场景}：
				                  \begin{itemize}[label={}, nosep]
					                  \item 客户端-服务器通信（C/S架构）
					                  \item 分布式系统通信
					                  \item 实时数据传输
				                  \end{itemize}
				            \item \textbf{代码示例}：
				                  \begin{lstlisting}[language=Java]
// Socket用于网络通信
Socket client = new Socket("127.0.0.1", 8080);
OutputStream out = client.getOutputStream();
out.write("Hello Server".getBytes());
client.close();
// Socket通过网络进行数据传输
				      \end{lstlisting}
			            \end{itemize}

			      \item[B] \textbf{ServerSocket用于服务器端 — 正确}
			            \begin{itemize}[label={}, leftmargin=*, nosep]
				            \item \textbf{功能说明}：监听指定端口，接受客户端连接
				            \item \textbf{工作流程}：
				                  \begin{itemize}[label={}, nosep]
					                  \item 绑定端口：new ServerSocket(port)
					                  \item 监听连接：serverSocket.accept()
					                  \item 获取客户端Socket：Socket client = accept()
					                  \item 使用客户端Socket通信
				                  \end{itemize}
				            \item \textbf{代码示例}：
				                  \begin{lstlisting}[language=Java]
import java.io.*;
import java.net.*;

public class EchoServer {
    public static void main(String[] args) {
        try {
            // 创建服务器端Socket，绑定端口8080
            ServerSocket serverSocket = new ServerSocket(8080);
            System.out.println("服务器启动，等待连接...");

            // 阻塞等待客户端连接
            Socket clientSocket = serverSocket.accept();
            System.out.println("客户端已连接");

            // 获取输入输出流
            BufferedReader in = new BufferedReader(
                new InputStreamReader(clientSocket.getInputStream()));
            PrintWriter out = new PrintWriter(
                clientSocket.getOutputStream(), true);

            // 读取客户端消息
            String inputLine;
            while ((inputLine = in.readLine()) != null) {
                System.out.println("收到：" + inputLine);
                out.println("Echo: " + inputLine);  // 回射给客户端
            }

            // 关闭资源
            in.close();
            out.close();
            clientSocket.close();
            serverSocket.close();
        } catch (IOException e) {
            e.printStackTrace();
        }
    }
}
// ServerSocket在服务器端监听端口并接受连接
				      \end{lstlisting}
			            \end{itemize}

			      \item[C] \textbf{Socket用于客户端 — 正确}
			            \begin{itemize}[label={}, leftmargin=*, nosep]
				            \item \textbf{功能说明}：主动连接服务器，进行数据传输
				            \item \textbf{工作流程}：
				                  \begin{itemize}[label={}, nosep]
					                  \item 创建Socket：new Socket(host, port)
					                  \item 获取输入输出流
					                  \item 读写数据
					                  \item 关闭连接
				                  \end{itemize}
				            \item \textbf{代码示例}：
				                  \begin{lstlisting}[language=Java]
import java.io.*;
import java.net.*;

public class EchoClient {
    public static void main(String[] args) {
        try {
            // 创建客户端Socket，连接服务器
            Socket socket = new Socket("127.0.0.1", 8080);
            System.out.println("已连接到服务器");

            // 获取输入输出流
            PrintWriter out = new PrintWriter(
                socket.getOutputStream(), true);
            BufferedReader in = new BufferedReader(
                new InputStreamReader(socket.getInputStream()));
            BufferedReader stdIn = new BufferedReader(
                new InputStreamReader(System.in));

            // 从控制台读取并发送给服务器
            String userInput;
            while ((userInput = stdIn.readLine()) != null) {
                out.println(userInput);  // 发送给服务器
                System.out.println("服务器响应：" + in.readLine());
            }

            // 关闭资源
            out.close();
            in.close();
            stdIn.close();
            socket.close();
        } catch (IOException e) {
            e.printStackTrace();
        }
    }
}
// Socket在客户端连接服务器并进行通信
				      \end{lstlisting}
			            \end{itemize}

			      \item[D] \textbf{不需要处理异常 — 错误（本题答案）}
			            \begin{itemize}[label={}, leftmargin=*, nosep]
				            \item \textbf{错误原因}：网络操作不稳定，必须处理异常
				            \item \textbf{网络异常来源}：
				                  \begin{itemize}[label={}, nosep]
					                  \item 网络未连通（ConnectException）
					                  \item 主机不可达（UnknownHostException）
					                  \item 连接超时（SocketTimeoutException）
					                  \item 连接断开（SocketException）
					                  \item IO错误（IOException）
				                  \end{itemize}
				            \item \textbf{必须处理的异常类型}：
				                  \begin{lstlisting}[language=Java]
// Socket编程常见异常
try {
    Socket socket = new Socket("host", port);
    // 可能导致异常的网络操作
} catch (UnknownHostException e) {
    // 主机名无法解析
    e.printStackTrace();
} catch (ConnectException e) {
    // 连接被拒绝（服务器未启动）
    e.printStackTrace();
} catch (SocketTimeoutException e) {
    // 连接超时
    e.printStackTrace();
} catch (IOException e) {
    // 其他IO异常
    e.printStackTrace();
}
// 网络编程必须处理各种可能发生的异常
					      \end{lstlisting}
				            \item \textbf{代码反证}：
				                  \begin{lstlisting}[language=Java]
// 错误示例：不处理异常（编译不通过）
public class BadSocketExample {
    public static void main(String[] args) {
        Socket socket = new Socket("127.0.0.1", 8080);
        // ^ 编译错误！必须要处理IOException
    }
}

//javac错误信息：
//error: 未报告的异常错误IOException; 必须对其进行捕获或声明以便抛出
//    Socket socket = new Socket("127.0.0.1", 8080);
//                                      ^
//1 个错误
// Java强制要求处理IOException及其子类异常
					      \end{lstlisting}
				            \item \textbf{正确处理方式1：try-catch}
				                  \begin{lstlisting}[language=Java]
public class CorrectSocketExample {
    public static void main(String[] args) {
        try (Socket socket = new Socket("127.0.0.1", 8080)) {
            // 使用socket通信
        } catch (UnknownHostException e) {
            System.out.println("主机名解析失败");
        } catch (IOException e) {
            System.out.println("连接失败或IO错误");
            e.printStackTrace();
        }
    }
}
// 使用try-with-resources自动关闭，捕获所有可能异常
					      \end{lstlisting}
				            \item \textbf{正确处理方式2：throws声明}
				                  \begin{lstlisting}[language=Java]
public class ThrowsSocketExample {
    public static void main(String[] args) throws IOException {
        Socket socket = new Socket("127.0.0.1", 8080);
        // 使用socket
        socket.close();
    }
}

// 或者
public class MethodExample {
    public void connectToServer(String host, int port) throws IOException {
        Socket socket = new Socket(host, port);
        // 使用socket
        socket.close();
    }
}
// 将异常抛出给调用者处理（不推荐main方法这样做）
					      \end{lstlisting}
				            \item \textbf{健壮的服务器代码}：
				                  \begin{lstlisting}[language=Java]
public class RobustServer {
    public static void main(String[] args) {
        ServerSocket serverSocket = null;
        try {
            serverSocket = new ServerSocket(8080);
            while (true) {
                Socket client = serverSocket.accept();
                new Thread(() -> {
                    try {
                        handleClient(client);
                    } catch (IOException e) {
                        System.out.println("客户端通信异常：" + e.getMessage());
                    } finally {
                        try {
                            client.close();
                        } catch (IOException e) {}
                    }
                }).start();
            }
        } catch (IOException e) {
            System.err.println("服务器启动失败：" + e.getMessage());
        } finally {
            if (serverSocket != null) {
                try {
                    serverSocket.close();
                } catch (IOException e) {}
            }
        }
    }

    private static void handleClient(Socket client) throws IOException {
        // 处理客户端请求
    }
}
// 每个网络操作都必须处理IOException
					      \end{lstlisting}
			            \end{itemize}
		      \end{itemize}

		      \textbf{【应背诵知识点】}
		      \begin{itemize}[label={}, nosep]
			      \item Socket用于网络通信
			      \item ServerSocket服务器端监听，Socket客户端连接
			      \item 必须处理IO异常（UnknownHostException、IOException等）
			      \item 使用try-with-resources或try-catch处理异常
			      \item 口诀：Socket网络通通信，Server监听Socket连，所有操作要异常，不写捕获编译难
		      \end{itemize}

		      \textbf{【命题趋势与实战策略】}
		      \textbf{考场秒杀}：看到"不需要处理异常"→错，网络编程必须处理IOException
		      \textbf{高频陷阱}：忘记处理IOException导致编译错误
	      \end{bluebox}

\end{enumerate}

\end{document}
